% ======================================================================
% 第 IV 节：强耦合近似
% ======================================================================
\section{强耦合近似}\label{sec:strong}

决定夸克是否被束缚的规范场平均 $Z(P)$ 已在第三节
（式~\eqref{eq:3.13}）中在格点上定义。这一平均有两个极限可以计算。最有趣的是
强耦合极限 $g \to \infty$。这一极限展示夸克束缚。本节将导出强耦合展开。强耦合
展开将作为把规范场理论改写为弦模型的基础。这在本节也会说明。

具体考虑式~\eqref{eq:3.13} 的分子，记为 $I_{\theta}(P)$：
\begin{equation}
I_{\theta}(P) = \prod_{n,\mu\nu}\int\ddth{\mu\nu}
  \exp\!\Bigl[i\sum_{P}(\pm)\thl{\mu\nu}
    +\frac{1}{2g^{2}}\sum_{n,\mu\nu}
      \bigl(e^{if_{\mu\nu}(n)}+e^{-if_{\mu\nu}(n)}\bigr)\Bigr].
\label{eq:4.1}
\end{equation}

把被积函数按 $1/g^{2}$ 的幂展开，零阶项为
\begin{equation}
I_{\theta}^{(0)}(P) = \prod_{n,\mu\nu}\int\ddth{\mu\nu}
  \exp\!\Bigl[i\sum_{P}(\pm)\thl{\mu\nu}\Bigr].
\label{eq:4.2}
\end{equation}
这一项为零，因为对 $\sum_{P}$ 中出现的任何 $\thl{\mu\nu}$，积分
$\int\ddth{\mu\nu}\exp(\pm i\thl{\mu\nu})$ 为零。因此必须寻找能抵消线积分中
$\thl{\mu\nu}$ 的更高阶的 $g^{-2}$ 项。一阶项为
\begin{equation}
I_{\theta}^{(1)}(P) = \prod_{n,\mu\nu}\int\ddth{\mu\nu}
  \exp\!\Bigl[i\sum_{P}(\pm)\thl{\mu\nu}+if_{\mu\nu}(l)\Bigr].
\label{eq:4.3}
\end{equation}
量 $f_{\mu\nu}(l)$ 本身就是规范场的线积分；它是围绕从格点 $l$ 出发的、大小为
$a$（单位方块）的小方块的线积分。除非能找到一个单位方块，使 $f_{\mu\nu}(l)$
完全抵消线积分 $\sum_{P}(\pm)\thl{\mu\nu}$，否则 $I_{\theta}^{(1)}(P)$ 的积分
为零。这只有当路径 $P$ 本身是一个单位方块时才可能。否则一阶项为零，必须研究
$g^{-2}$ 或更高阶的项。

$g^{-2n}$ 阶的项具有形式
\begin{equation}
I_{\theta}^{(n)}(P) = \prod_{n,\mu\nu}\int\ddth{\mu\nu}
  \exp\!\Bigl[i\sum_{P}(\pm)\thl{\mu\nu}
    + i\sum_{m=1}^{n}f_{\mu\nu}(l_{m})\Bigr].
\label{eq:4.4}
\end{equation}
这些求和中非零的项只是那些满足
\begin{equation}
\sum_{m}f_{\mu\nu}(l_{m}) = 0.
\label{eq:4.5}
\end{equation}
[关于用 $\thl{\mu\nu}$ 表示 $f_{\mu\nu}$ 的定义见式~\eqref{eq:3.8}。] 这一方程
可以从几何上理解。每个 $f_{\mu\nu}$ 对应格点上一个大小为 $a$ 的小方块。要使该和
为零，由 $f_{\mu\nu}$、$f_{\mu\nu'}$、…… 定义的小方块组必须组合成以 $P$ 为边界
的曲面。（确切地说，每个 $f_{\mu\nu}$ 对应绕一个小方块的线积分，当这些方块拼成
曲面时，线积分必须沿曲面所有内部线抵消。曲面边界 $P$ 上的线积分必须与原来的
路径 $P$ 反向。）见图~\ref{fig:10}。

对给定的路径 $P$，$I_{\theta}(P)$ 中最低的非零阶由 $P$ 所围的最小面积 $A$ 决定，
面积 $A$ 是格点上以 $P$ 为边界、由单位方块构成的任意曲面的面积。于是，除数值因子
外，$I_{\theta}(P) \sim (g^{-2})^{A/a^{2}}$。

\begin{figure}[t]
  \centering
  \includegraphics[width=0.5\textwidth]{fig10}
  \refstepcounter{figure}%
  \caption*{图 10. 用单位方块填满路径 $P$ 所围的面积。}
  \label{fig:10}
\end{figure}

这就是第二节所承诺的结果：规范场平均 $I_{\theta}(P)$ 的行为如同
$\exp[-A\,(\ln g^{2})/a^{2}]$，即对 $P$ 所围面积呈指数形式。因此，根据第二节的
论证，夸克路径不会宏观地分离，末态粒子中将没有夸克。

考虑固定 $P$ 时 $I_{\theta}(P)$ 展开中的高阶项。这样的项很多，因为以 $P$ 为边界
的曲面很多。特别是，组合 $f$ 的子集使之相加为零的方式很多，因此这样的子集可以
加到任何构成以 $P$ 为边界的曲面的 $f$ 最小和上。相加为零的 $f$ 集合的最简单例子
是对应于单位立方体六个面的 $f$ 集合。把它写出来，得到
\begin{equation}
\begin{split}
&f_{\mu\nu}(n)+f_{\mu\nu}(n+\hat\mu)+f_{\mu\nu}(n+\hat\nu)\\
&\qquad + f_{\mu\rho}(n)+f_{\rho\mu}(n+\hat\mu)+f_{\mu\rho}(n+\hat\rho)=0,
\end{split}
\label{eq:4.6}
\end{equation}
利用式~\eqref{eq:3.8} 容易验证。[式~\eqref{eq:4.6} 是方程
$\epsilon_{\mu\nu\rho\sigma}\partial_{\rho}F_{\mu\nu}(x)=0$ 的格点模拟。]

设 $A(P)$ 是上面定义过的、由 $P$ 所围的最小面积。由于可以在这一最小面积上任意
放置单位立方体，这意味着在 $I_{\theta}(P)$ 的展开中，$g^{-2(n+1)}$ 阶的项大约
比 $g^{-2n}$ 阶的项多 $A(P)/a^{2}$ 个。[可以在空间中任意放置单位立方体，而不
仅限于最小曲面；但当用 $Z$ 除 $I_{\theta}(P)$ 时，所有不相连的项都如常抵消。]
这表明 $1/g^{2}$ 展开在 $A(P) \to \infty$ 极限下没有多大用处，而这一极限正是夸克
束缚所关心的极限。不过，处理相关问题时的经验表明，$I_{\theta}(P)$ 并不是应该
展开的恰当量；相反，应当尝试写成
\[
I(P) = Z^{-1}I_{\theta}(P) = \bigl(g^{-2}\bigr)^{A(P)/a^{2}}e^{-c(P,g^{2})}
\]
并代之以按 $g^{-2}$ 的幂展开 $c(P,g^{2})$。可以预期 $c(P,g^{2})$ 由正比于 $A(P)$
的项主导，比如
\[
c(P,g^{2}) = A(P)f(g^{-2}) + \order{P}
\]
（其中 $P$ 是路径 $P$ 的长度）。关键问题是 $f(g^{-2})$ 级数的性质。处理类似类型的
展开（即统计力学的高温展开：例如见文献 13）的过往经验表明，至少对小于临界值
$g_{c}^{-2}$ 的 $g^{-2}$，$f(g^{-2})$ 会有收敛展开。不过，目前还没有人在规范场理论
中计算过 $f(g^{-2})$。

考虑 $I(P)$ 的完整展开。展开中每一项对应一个周长为 $P$ 的曲面。完整展开对应
对给定周长 $P$ 的所有可能曲面求和。``所有可能''曲面包括自交曲面（以计入给定
$f_{\mu\nu}$ 在和 $\sum f_{\mu\nu} + f_{\mu\nu'} + \cdots$ 中出现多次的项）。
除由曲面面积确定的 $g^{-2}$ 幂次外，每个曲面还有一个权重因子。对简单曲面，权重
因子为 1；对自交曲面权重更复杂。

因此，流-流传播子的强耦合展开与强子弦模型具有相同的总体结构。这里实际上处理
的是二重展开。狄拉克场作用量中出现的系数 $\kappa$ 的展开用来定义格点上的夸克圈；
$\kappa$ 展开之和是对所有可能夸克圈的求和。$g^{-2}$ 展开用来定义填充夸克圈的
曲面。$g$ 展开之和是对所有这类曲面的求和。这正是弦模型中出现的结构：夸克圈与
插值曲面组合求和。不过，规范场理论的圈与曲面定义在格点上，而弦模型的圈与曲面
定义在连续体上。要导出这两类曲面之间的定量关系可能并不容易。
